% ============================================================
% The Coherence Decomposition
% ============================================================

\documentclass[12pt]{amsart}

\usepackage{amsmath,amssymb,amsthm}
\usepackage{booktabs}
\usepackage{microtype}
\usepackage{url}

\newtheorem{theorem}{Theorem}
\newtheorem{lemma}[theorem]{Lemma}
\newtheorem{proposition}[theorem]{Proposition}
\newtheorem{corollary}[theorem]{Corollary}
\newtheorem{conjecture}[theorem]{Conjecture}
\theoremstyle{definition}
\newtheorem{definition}[theorem]{Definition}
\theoremstyle{remark}
\newtheorem{remark}[theorem]{Remark}

\title{The Coherence Decomposition}

\author{Alexander S.\ Petty}

\email{alexander.petty@gmail.com}
\date{June 2022 (revised April 2026)}
\subjclass[2020]{11A63, 11B83}

\begin{document}
\begin{abstract}
We introduce two new invariants of the fractional field
$\mathcal{F}(n) = \{k/n : 1 \le k \le n{-}1\}$: the \emph{pairwise
alignment} $\sigma(n)$ and the \emph{focused alignment}
$F(n) = \alpha(n) - \sigma(n)$, where $\alpha(n)$ is the repetend
alignment of previous papers. For digit-partitioning primes, these
admit exact formulas:
\[
\alpha(ps) = \frac{2s-1}{ps-1}, \qquad
F(ps) = \frac{s}{ps-1}, \qquad
\sigma(ps) = \frac{s-1}{ps-1},
\]
with limits $2/p$, $1/p$, and $1/p$ respectively as the smooth
factor $s \to \infty$.

The decomposition $\alpha = F + \sigma$ separates total coherence into
a focused component (how strongly $1/n$ is distinguished within its
field) and a diffuse component (the background pairwise coherence
among all fractions). For digit-partitioning primes, the golden threshold
$1/\varphi$ lies strictly between the focused and total
limits if and only if $p \in \{2, 3\}$. Among primes
$p \ge 3$, only $p = 3$ has this property.
\end{abstract}
\maketitle

% ============================================================
\section{Introduction}

This paper decomposes alignment into focused and diffuse components.

The preceding papers in this series studied the repetend alignment
$\alpha(n)$, which measures the coherence of the fractional field
$\mathcal{F}(n)$ relative to $1/n$. This paper introduces a
complementary measure: the \emph{pairwise alignment} $\sigma(n)$,
which quantifies the background coherence among all pairs of
fractions in the field, independent of any distinguished member.

The difference $F(n) = \alpha(n) - \sigma(n)$ isolates the
\emph{focused alignment}: the excess coherence that $1/n$ enjoys
over a generic field member. Together, $\alpha$, $\sigma$, and $F$
decompose the coherence of $\mathcal{F}(n)$ into structural
components with distinct limits and distinct roles in the golden
ratio classification.

% ============================================================
\section{Definitions}

\begin{definition}
The \emph{pairwise alignment} of $n$ in base $b$ is
For $n \ge 3$:
\[
\sigma(n) = \binom{n-1}{2}^{-1}
\sum_{1 \le j < k \le n-1}
a(j/n,\, k/n).
\]
where $a(j/n,\, k/n)$ is the position-wise digit-match proportion
between the repetends of $j/n$ and $k/n$. Pairs of terminating
fractions contribute $a = 1$. Pairs of non-terminating fractions
with the same cycle length are compared position-wise. All other
pairs (terminating with non-terminating, or different cycle
lengths) contribute $a = 0$.
\end{definition}

\begin{definition}
The \emph{focused alignment} of $n$ is
\[
F(n) = \alpha(n) - \sigma(n).
\]
\end{definition}

Focused alignment measures how much more coherent the field is with
$1/n$ than with a randomly chosen pair of its own members.

% ============================================================
\section{The Decomposition for Digit-Partitioning Primes}

\begin{theorem}\label{thm:decomposition}
Let $p$ be a digit-partitioning prime in base $b$ and let $s \ge 1$
be a $b$-smooth integer. Then
\begin{align}
\alpha(ps) &= \frac{2s-1}{ps-1}, \label{eq:alpha} \\[4pt]
F(ps) &= \frac{s}{ps-1}, \label{eq:F} \\[4pt]
\sigma(ps) &= \frac{s-1}{ps-1}. \label{eq:sigma}
\end{align}
\end{theorem}

\begin{proof}
Equation~\eqref{eq:alpha} is the alignment formula from~\cite{paper2}.

For $\sigma$, we count the pairs $(j,k)$ with $j < k$ and
$a(j/n, k/n) = 1$. There are three sources of such pairs.

\smallskip
\emph{Terminating--terminating pairs.}\quad
The $s - 1$ terminating fractions (multiples of $p$ in
$\{1, \ldots, ps{-}1\}$) all contribute $a = 1$ in pairs by the
definition of pairwise alignment, giving $\binom{s-1}{2}$ pairs.

\smallskip
\emph{Identical-repetend pairs among non-terminating fractions.}\quad
By the cycle-partition argument in~\cite{paper2}
(Universal Alignment Formula section),
the $(p-1)s$ non-terminating fractions in $\{k/(ps)\}$ partition
into $(p-1)/L$ families of $sL$ each, with $L = \operatorname{ord}_p(b)$.
Within each family, the fractions split into $L$ same-repetend
classes of $s$ each, indexed by the cyclic shift of the base
repetend. (Each of the $L$ shifts is achieved by exactly $s$
fractions, since the residue $r \in \{1, \ldots, p{-}1\}$
producing a given shift has $s$ representatives in
$\{k : 1 \le k \le ps{-}1,\ p \nmid k\}$.) Across all
$(p-1)/L$ families, there are $(p-1)/L \cdot L = p - 1$
same-repetend classes total, each of size $s$. Within each
same-repetend class, all $\binom{s}{2}$ pairs have identical
repetends and contribute $a = 1$. The total is
$(p-1)\binom{s}{2}$ pairs.

\smallskip
\emph{All other pairs contribute $0$.}\quad
Pairs of non-terminating fractions with different repetends
(i.e., from different same-repetend classes) are cyclic shifts
of each other; by the digit-partitioning characterization
in~\cite{paper2}, the digit function is injective on
$\{1, \ldots, p{-}1\}$, so distinct cyclic shifts share no
positions and contribute $a = 0$. Pairs mixing a terminating
with a non-terminating fraction contribute $a = 0$ by definition.

\smallskip
Summing:
\[
\Sigma = \binom{s-1}{2} + (p-1)\binom{s}{2}.
\]

Expanding and dividing by $\binom{ps-1}{2}$:
\begin{align*}
\sigma(ps)
&= \frac{\binom{s-1}{2} + (p-1)\binom{s}{2}}{\binom{ps-1}{2}}
= \frac{(s-1)(s-2)/2 + (p-1)s(s-1)/2}{(ps-1)(ps-2)/2} \\
&= \frac{(s-1)\bigl[(s-2) + (p-1)s\bigr]}{(ps-1)(ps-2)}
= \frac{(s-1)(ps - 2)}{(ps-1)(ps-2)}
= \frac{s-1}{ps-1}.
\end{align*}

Finally, $F = \alpha - \sigma = (2s{-}1)/(ps{-}1) - (s{-}1)/(ps{-}1)
= s/(ps{-}1)$.
\end{proof}

\begin{corollary}\label{cor:limits}
As $s \to \infty$ through $b$-smooth integers,
\[
\alpha(ps) \to \frac{2}{p}, \qquad
F(ps) \to \frac{1}{p}, \qquad
\sigma(ps) \to \frac{1}{p}.
\]
\end{corollary}

At infinite resolution, total coherence splits equally between
focused and diffuse components.

% ============================================================
\section{The Golden Gap}

\begin{theorem}\label{thm:golden-gap}
For digit-partitioning primes $p$ not dividing $b$, the
golden threshold $1/\varphi$ lies strictly between the focused
and total coherence limits, $1/p < 1/\varphi < 2/p$, if and only
if $p \in \{2, 3\}$. For $p = 2$, the digit-partitioning
formula gives $\alpha(2s) = 1$ identically, so the
threshold question is trivial. Among primes $p \ge 3$,
only $p = 3$ has the gap.
\end{theorem}

\begin{proof}
The condition $1/p < 1/\varphi$ requires $p > \varphi \approx 1.618$,
satisfied by all primes. The condition $1/\varphi < 2/p$ requires
$p < 2\varphi \approx 3.236$, satisfied only by $p = 2$
and $p = 3$.
\end{proof}

For $p = 3$, the limits are $F \to 1/3 \approx 0.333$ and
$\alpha \to 2/3 \approx 0.667$, with $1/\varphi \approx 0.618$
falling in between. The golden threshold occupies the gap between
focused and total coherence. For $p \ge 5$, the total limit $2/p
\le 2/5 = 0.4$ falls below $1/\varphi$, so both $F$ and $\alpha$
remain below the golden threshold at all resolutions.

This gives a new characterization of the golden ratio's role. The
three-tier classification arises because $p = 3$ is the unique prime
whose diffuse coherence $\sigma$ is large enough to push the total
$\alpha = F + \sigma$ above $1/\varphi$. The focused component alone
never reaches $1/\varphi$ for any prime. It is the accumulation of
pairwise background coherence that carries the total across the
threshold.

% ============================================================
\section{The Sigma-Zero Characterization}

\begin{proposition}\label{prop:sigma-zero}
For a prime $p \ge 3$ not dividing $b$:
$\sigma(p) = 0$ if and only if $p$ is
digit-partitioning in base $b$.
\end{proposition}

\begin{proof}
If $p$ is digit-partitioning, no two distinct fractions
$j/p$ and $k/p$ share a digit at any common position,
so every off-diagonal $a(j/p, k/p) = 0$ and
$\sigma(p) = 0$.

Conversely, if $p$ is not digit-partitioning
($p > b{+}1$), there exist distinct $r, s$ with
$\delta(r) = \delta(s)$. The fractions $r/p$ and $s/p$
share the digit $\delta(r)$ at position $0$, giving
$a(r/p, s/p) \ge 1/L > 0$ where
$L = \operatorname{ord}_p(b)$. Hence $\sigma(p) > 0$.
\end{proof}

This provides an equivalent characterization of the
digit-partitioning property: $p$ is digit-partitioning if and
only if no two
distinct fractions in $\{k/p : 1 \le k \le p{-}1\}$ share a digit
at any common position.

% ============================================================
\section{Composite Denominators}

For composite rough parts, the decomposition
$\alpha = F + \sigma$ continues to hold by definition,
but the exact formulas for digit-partitioning primes
do not extend directly. The composite case requires a
theory of cross-cycle digit coincidences that is beyond
the scope of this paper.

% ============================================================
\section{Remarks}

\subsection*{Three invariants}

The fractional field $\mathcal{F}(n)$ now carries three defined
invariants:
\begin{itemize}
\item $\alpha(n)$: total coherence (alignment with $1/n$),
\item $\sigma(n)$: diffuse coherence (pairwise background alignment),
\item $F(n)$: focused coherence (structural identity of $1/n$).
\end{itemize}
These are distinct measures with distinct limits and distinct roles
in the classification. Alpha connects to the golden ratio. Sigma
characterizes digit-partitioning at prime level. $F$ measures
the excess coherence of $1/n$ over a generic pair. All
three are computable from the fractional field alone.

\subsection*{The gap as a structural principle}

The golden threshold does not act on a single quantity. It acts on
the \emph{gap} between focused and total coherence. The three-tier
classification arises because $p = 3$ is the unique prime for which
this gap is wide enough to contain $1/\varphi$. The width of the gap
is $1/p$ (the diffuse limit), and the condition $1/p > 1/\varphi -
1/p$ simplifies to $2/p > 1/\varphi$, recovering the golden bound
$p < 2\varphi$.

% ============================================================
\begin{thebibliography}{9}

\bibitem{paper2}
A.~S.~Petty,
Digit-partitioning primes and the alignment formula,
research note, 2020.

\bibitem{hardy}
G.~H.~Hardy and E.~M.~Wright,
\emph{An Introduction to the Theory of Numbers},
6th ed., Oxford University Press, 2008.

\bibitem{davenport}
H.~Davenport,
\emph{Multiplicative Number Theory},
3rd ed., Springer, 2000.

\bibitem{montgomery}
H.~L.~Montgomery and R.~C.~Vaughan,
\emph{Multiplicative Number Theory I: Classical Theory},
Cambridge University Press, 2007.

\end{thebibliography}

\end{document}
